\documentclass[11pt]{article}
\usepackage{amsmath, amssymb}
\usepackage{hyperref}
\usepackage{geometry}
\geometry{a4paper, margin=1in}
\usepackage{graphicx} 
\usepackage{float}
\usepackage{setspace}
\onehalfspacing

\title{Unitary Phase Representations in Recurrent Neural Networks: A Path to Modular and Reversible Intelligence}
\author{Matthew Murphy}
\date{February 22, 2026}

\begin{document}

\maketitle

\begin{abstract}
The advancement of Recurrent Neural Networks (RNNs) has been structurally bottlenecked by the Contraction Problem: a mathematical necessity for standard gate-based models, such as LSTMs and GRUs, to dissipative their internal state to maintain numerical stability. While this prevents explosive divergence, it induces a chronic decay of information fidelity over time, rendering standard intelligence lossy and irreversible. This paper presents the Riemannian Optimized Unified Neural Dynamo (ROUND), an architecture that replaces contractive dissipation with Unitary Phasic Resonance. By representing state as a position on a high-dimensional Riemannian manifold equipped with a metric tensor ($g$), ROUND achieves Unitary Isometry, which perfectly preserves the distance between representations ($\|P(x)\| = \|x\|$), as well as Topological Isomorphism, where the model structural matches the complex geometry of the processed signal. Through five targeted benchmarks, we demonstrate that ROUND is an operationally superior substrate for modular, reversible, and high-fidelity machine intelligence.
\end{abstract}

\newpage

\section{The Contraction Problem in Classical RNNs}

In standard recurrent architectures, stability is achieved through Contraction Analysis [1], which requires that a state-update function $f$ be a contractive mapping ($|f(x) - f(y)| < |x - y|$). Traditionally, this is implemented via nonlinear gating mechanisms, such as sigmoid or tanh functions, that systematically "squash" inputs toward a zero-point attractor. The consequence of this approach is incremental Signal Decay: information is continuously erased to prevent gradient explosion, creating a "vanishing gradient fog" where long-term temporal dependencies are lost. This process is fundamentally irreversible; once a signal is collapsed through a contractive gate, the original input cannot be recovered algebraically. To the standard RNN, learning is a stochastic search through a fading echo, limited by the lossy nature of its internal state-space compression.

\section{Architectural Innovation: Orthogonal Phase Topology}

The ROUND architecture represents a fundamental departure from the contractive paradigm by mapping internal state transitions to a high-dimensional Riemannian Manifold. Specifically, the model utilizes a hypertorus topology where curvature and distance are governed by Unitary Isometry [2].

\subsection{Core Mathematical Definitions}

\textbf{Unitary Isometry:} A distance-preserving mapping where the inner product (and thus the norm) of the state vector remains invariant ($\|P(x)\| = \|x\|$). In the ROUND neuron, the hidden state’s magnitude is identically $1.0$, ensuring that informatic density is neither diluted nor amplified across recurrent cycles. This shift from linear displacement to unitary rotation ensures that stored information is resistant to dissipative decay.

\textbf{Topological Isomorphism:} A mapping that preserves the underlying topological structure between the model’s internal phase-space and the geometric manifold of the task. This enables a phenomenon known as Search Collapse, where the network "locks" onto resonance with the signal geometry rather than iteratively approximating it through traditional error reduction.

\textbf{Uniform Phase Topology (UPT):} The structural requirement that data be represented via phase angles ($\phi$) distributed uniformly across the manifold [6]. This prevents the "informatic blotting" common in standard RNNs, ensuring that every region of the latent space remains equally accessible and resistant to gravitational collapse toward the origin.

\textbf{Riemannian Manifold:} A smooth mathematical space ($M$) equipped with a metric tensor ($g$) that allows for the calculation of curvature-aware geodesics [3]. ROUND exists on this manifold, navigating state transitions via unitary rotations ($h_{t} = R \cdot h_{t-1}$) where the internal updates are expressed in $U = x + \epsilon i$ coordinates, decoupling classical observables from infinitesimal rotational states.

\newpage

\section{Validated Capabilities and Benchmarks}

The theoretical claims of the ROUND architecture are quantified through five experimental benchmarks, each designed to expose a fundamental failure mode of dissipative recurrent systems.

\subsection{1. The Crystalline Loop: Recursive Stability}

\textbf{Method:} A deep recurrent stack processes 8-bit ASCII bit-streams through a continuous recurrent loop to measure long-term persistence.

\begin{figure}[H]
    \centering
    \includegraphics[width=0.75\linewidth]{verification_report_2026-01-27_0140_SCIENTIFIC_PROTOCOL_TARGETED.png}
    \caption{\textit{Verification of Topological Isomorphism.}\textit{ Comparison of bitwise retrieval accuracy. By mapping state to a high-dimensional Riemannian manifold, ROUND preserves the geometric distance between representations ($|P(x)| = |x|$), achieving perfect persistence where traditional contractive models dissolve into "vanishing gradient fog."}}
    \label{fig:1}
\end{figure}

\textbf{Findings:} While contractive GRUs exhibit stochastic drift and eventual collapse to the manifold's origin, ROUND maintains 100\% bitwise recovery across arbitrary sequence lengths. This confirms that Phasic Resilience allows for the indefinite preservation of symbolic identity within stable attractor basins.

\newpage

\subsection{2. The Relay Test: Zero-Shot Modularity}

\textbf{Method:} An Encoder and Decoder are trained in isolation and then coupled "zero-shot" to evaluate representational consistency.

\begin{figure}[H]
    \centering
    \includegraphics[width=0.75\linewidth]{sandwich_duel_story_2026-01-27_0140_SCIENTIFIC_PROTOCOL_TARGETED.png}
    \caption{\textit{The Sandwich Duel: Zero-Shot Associative Relay.} A bitstream is processed by an isolated decoder into a latent state, which is then relayed to an independently trained encoder for bitwise reconstruction. While standard models fail to synchronize their stochastic hidden spaces (0.4\% success), ROUND achieves perfect 100\% relay integrity. This confirms that Topological Isomorphism creates a universal, modular language across disparate systems.}
    \label{fig:2}
\end{figure}

\textbf{Findings:} The "Sandwich Duel" confirms that ROUND achieves 100\% reconstruction integrity in zero-shot configurations. This result demonstrates that the underlying Phasic Resonance creates a robust, invariant coordinate system that is consistent across independent training runs.

\newpage

\subsection{3. The Prism Stack: Geometric Logic}

\textbf{Method:} A modular addition task ($[L + P] \pmod{18}$) where a "Lens" input sets the initial phase state of the manifold.

\begin{figure}[H]
    \centering
    \includegraphics[width=0.75\linewidth]{prism_stack_duel_2026-01-27_0140_SCIENTIFIC_PROTOCOL_TARGETED.png}
    \caption{The Prism Stack: Phase-State Initialization and Resonant Logic. ROUND performs modular addition as a native geometric operation by mapping inputs to its rotational manifold. While standard RNNs must approximate modular boundaries through search, ROUND achieves 100\% accuracy by processing the "Light" input as a simple phase displacement relative to the "Lens" initialization. }
    \label{fig:3}
\end{figure}

\textbf{Findings:} ROUND achieves 100\% accuracy with near-instantaneous convergence because modular arithmetic is a native property of its internal rotational physics.

\newpage

\subsection{4. The Color Algebra Duel: Free Reversibility}

\textbf{Method:} The architecture characterizes the additive logic of discrete color mixtures within its latent phase-space.

\begin{figure}[H]
    \centering
    \includegraphics[width=0.75\linewidth]{color_algebra_duel_2026-01-27_0140_SCIENTIFIC_PROTOCOL_TARGETED.png}
    \caption{Color Algebra: Native Phase Arithmetic on the Riemannian Manifold. The network receives two discrete color indices and outputs their geometric midpoint. Because ROUND's hidden state inhabits a toroidal topology, it computes the midpoint as a native rotation, achieving 100\% accuracy. The ability to perform "phase arithmetic" natively confirms that ROUND's topology is algebraically functional—a necessary precondition for reversible, composable intelligence.}
    \label{fig:4}
\end{figure}

\textbf{Findings:} Because ROUND logic is expressed as unitary rotations, the operation is trivially reversible ($R^{-1} = R^{\top}$). The network acquired the inverse operation with zero additional training, demonstrating that reversible intelligence can reduce the computational overhead of bidirectional tasks by 50\% [4].

\newpage

\subsection{5. Riemannian Topological Recovery}

\textbf{Method:} High-fidelity reconstruction of high-frequency continuous signals (Sine Waves).

\begin{figure}[H]
    \centering
    \includegraphics[width=0.75\linewidth]{riemannian_recovery_2026-01-27_1130_SCIENTIFIC_PROTOCOL_TARGETED.png}
    \caption{Riemannian Topological Recovery: Curvature Tracking on a Quantized Manifold. The network must predict sinusoidal points by tracking its internal phase. Because ROUND's hidden state inhabits a toroidal phase manifold, it possesses a native geometric affinity for smooth curvature. The result is a high-precision reconstruction (MSE < 0.0001), confirming that a Riemannian topology is the natural substrate for representing continuous dynamics even when discretized for digital computation.}
    \label{fig:5}
\end{figure}

\textbf{Findings:} ROUND reached the "Crystalline Lock" threshold (MSE < 0.0001) with significantly higher precision than standard models, proving that Isomorphism is an operationally superior substrate for signal recovery.

\newpage

\section{Conclusion}

The benchmarks presented here confirm that the ROUND architecture is an operationally superior substrate for high-fidelity intelligence. By enforcing Unitary Phasic Topology, we enable a transition from opaque, lossy approximations to transparent, composable, and reversible systems. Furthermore, the architecture’s inherent composability suggests a path toward arbitrary scalability; because the core logic can be effectively quantized into a binary state-space, the computational overhead remains bounded even as model depth increases. This quantization facilitates the assembly of large-scale, modular systems where each neuron acts as a discrete, reversible processing unit. Ultimately, ROUND’s ability to achieve perfect isometry represents a fundamental shift in machine learning: moving from "learning by guessing" toward "intelligence by resonance."

\section{Selected References}

\begin{enumerate}
    \item \textbf{Lohmiller, W., \& Slotine, J. J. (1998).} "On Contraction Analysis for Non-linear Systems." *Automatica*.
    \item \textbf{Arjovsky, M., Shah, A., \& Bengio, Y. (2016).} "Unitary Recurrent Neural Networks." *ICML*.
    \item \textbf{Bernstein, E., \& Vazirani, U. (1997).} "Quantum Complexity Theory." *SIAM Journal on Computing*.
    \item \textbf{Landauer, R. (1961).} "Irreversibility and Heat Generation in the Computing Process." *IBM Journal of Research and Development*.
    \item \textbf{Vazza, F., \& Feletti, A. (2020).} "The Quantitative Comparison between the Neuronal Network and the Cosmic Web." *Frontiers in Physics*.
    \item \textbf{Naitzat, G., Zhitlukhin, A., \& Lekala, L. (2020).} "Topology of Deep Neural Networks." *International Conference on Machine Learning (ICML)*.
\end{enumerate}

\end{document}
